\section{Open Questions}

The replication reproduces the paper's stated PoC. The following five
questions are genuinely open --- they are neither answered by the paper
nor by this replication --- and each is scoped so a follow-up study
could actually attempt it.

\begin{enumerate}
\item \textbf{Noise-aware reward.}
Does adding a noise-aware reward term (e.g.\ minimize expected error under a
realistic depolarizing / $T_1$-$T_2$ noise model) change which schedule the
PPO agent converges to, and does it beat noise-oblivious ALAP on fidelity
rather than cycle count?
\emph{Basis:} qgym's PoC optimizes cycle count only; real compilation trades
depth for fidelity. \texttt{qgym.Rewarder} is pluggable but noise-awareness
is not exercised in the paper.
\emph{Next steps:} subclass \texttt{BasicRewarder} to add a per-gate fidelity
penalty from a Qiskit \texttt{NoiseModel} ($T_1$/$T_2$ + depolarizing);
retrain PPO $10^5$ steps at seed 20260703; evaluate PPO vs ALAP under the
same noise model on mean process fidelity and cycle count. Success $=$ PPO
strictly beats ALAP on fidelity even if it ties/loses on cycles.

\item \textbf{Modern RL algorithms.}
Do off-policy or model-based RL algorithms (SAC-Discrete, DreamerV3,
MuZero-style) reach ALAP parity in materially fewer environment steps than
vanilla PPO on the qgym Scheduling PoC, and do they scale where PPO stalls?
\emph{Basis:} the paper uses vanilla PPO only. Our replication showed PPO's
convergence is hyperparameter-sensitive. Off-policy replay + world-model
methods are typically far more sample-efficient on discrete combinatorial
tasks.
\emph{Next steps:} adapt qgym Scheduling to Gymnasium+cleanrl; train
SAC-Discrete and small DreamerV3 with matched seed on 3-qubit,
\texttt{max\_gates=5}. Compare (i) steps to first ALAP-parity, (ii) final
mean cycles on 50-circuit held-out set (seeds 20261703..20261752).

\item \textbf{Real hardware topologies.}
How does a qgym-trained Scheduling agent perform on real hardware coupling
maps (IBM Heron/Eagle heavy-hex, IonQ linear, Rigetti Aspen) as opposed to
the synthetic 3-qubit \texttt{MachineProperties} in the paper, and does the
trained policy transfer or require per-topology retraining?
\emph{Basis:} the PoC uses an abstract 3-qubit machine with no coupling
constraints. Whether the qgym RL approach generalizes to production hardware
is untested.
\emph{Next steps:} instantiate three qgym \texttt{MachineProperties} matching
(i) IBM Eagle 27q heavy-hex sub-lattice, (ii) IonQ 11q linear, (iii) Rigetti
Aspen-M fragment. Train PPO per topology on random 10-gate circuits, then
cross-evaluate. Metric: zero-shot mean-cycle delta vs each topology's ALAP.
Success $=$ within 10\% of ALAP zero-shot on unseen topologies.

\item \textbf{Transfer learning across topologies.}
Can a policy pre-trained on topology $A$ be fine-tuned to topology $B$ in
far fewer steps than training from scratch, and does a graph-neural-network
state encoder over the coupling graph beat qgym's default MLP for transfer?
\emph{Basis:} transfer across coupling maps is the practically useful version
of question~3. qgym's default \texttt{MultiInputPolicy} uses a flat MLP over
dict observations; GNNs are the standard architectural choice for graph-RL.
\emph{Next steps:} implement a GNN feature extractor (nodes $=$ qubits,
edges $=$ coupling). Pre-train PPO on topology $A$ for $10^5$ steps; continue
on $B$ for $\{10^4, 5{\times}10^3, 2.5{\times}10^3\}$ steps. Compare to
training from scratch on $B$ at same budgets. Metric: cycles-to-parity vs
cumulative environment steps.

\item \textbf{ML-in-the-loop with a classical compiler.}
In a hybrid setup, can a qgym-trained agent be inserted as a heuristic inside
Qiskit's \texttt{PassManager} (or a \texttt{t|ket>} pass) to strictly improve
on the classical compiler's schedule depth on realistic (20--50 qubit)
benchmark circuits, or does the classical compiler dominate at that scale?
\emph{Basis:} the most useful home for a qgym agent is not to \emph{replace}
a classical compiler but to be a callable subroutine inside one. The paper
never tests such an integration.
\emph{Next steps:} write a Qiskit \texttt{TransformationPass} that calls a
trained qgym Scheduling agent to schedule each basic block, keeping Qiskit's
routing/synthesis passes unchanged. Run hybrid vs stock
\texttt{transpile(optimization\_level=3)} on QASMBench medium (20q) and
QASMBench large (40--50q). Compare depth, two-qubit gate count, transpile
wall time. Success $=$ hybrid strictly beats stock Qiskit on depth without
regressing gate count.
\end{enumerate}
